% !PRINTED 165
\weberpage{165}
%<<CONT
is divisible by $w$, and that, if $U$, $U'$, $U''\ldots$ are integral rational
functions divisible by $w$, and $V$, $V'$, $V''\ldots$ are arbitrary integral
rational functions, then $UV + U'V' + U''V''\ldots$ is also divisible by $w$.

If $W$ is divisible by $w$, one also says that $w$ goes into $W$.

Two integral rational functions $U$, $V$ that are not divisible by one and the
same integral rational function are called \gesperrt{relatively prime}
or \gesperrt{coprime}.

We prove the theorem:

I. \gesperrt{If $U$, $V$, $v$ are integral rational functions of any
variables whatever, if $U$ and $v$ are relatively prime and $UV$ is divisible
by $v$, then $V$ is divisible by $v$.}

This theorem corresponds exactly to a well-known fundamental theorem from the
theory of the integers, namely that if a product of two integers is divisible
by a third integer which is coprime to one of the factors, then the other
factor must be divisible by this number.

We prove it by complete induction. If $U$, $V$, $v$ depend on only
\gesperrt{one} variable $x$, the theorem is true; for by \S.\,6 (Theorem I)
one can in this case, when $U$ and $v$ are relatively prime, determine two
other integral functions $P$ and $p$ of $x$ so that
\[
PU + pv = 1,
\]
from which, by multiplication by $V$,
\[
PUV + pVv = V
\]
follows, and from this one sees that if $UV$ is divisible by $v$, then $V$
must also be divisible by $v$.

We therefore assume that the theorem we wish to prove has been proved for
functions of $n$ and fewer variables, and from this we derive its validity
for functions of $n + 1$ variables.

For this it is necessary that we draw from Theorem I, assumed to be true,
some consequences, whose conclusion then yields the validity of the theorem
for the next higher number of variables.

If $v$ is an indecomposable function, then another function $U$ of the same
variables is either relatively prime to $v$ or divisible by $v$. From this it
follows by I. that a product of
%CONT>>
